\subsection{等变函数}

在本小节中，我们取定 G 的一个闭子群 H。用 \(\varpi\) 表示 G 到 G/H 上的典范投射。

除 G 上的哈尔测度 \(\mu\) 外，我们再取定 H 上的一个左哈尔测度 \(\beta\)。

根据 INT, VII, 第 56 页, § 2, 第 5 目, 定理 2，存在连续函数 \(\kappa\colon\mathbf{G}\to\mathbf{R}_{+}^{*}\)，使得 \(\kappa(xh)=\Delta_{\mathrm{H}}(h)\Delta_{\mathrm{G}}(h)^{-1}\kappa(x)\) 对每个 \((x,h)\in\mathbf{G}\times\mathbf{H}\) 成立。我们取定这样一个函数 \(\kappa\)。用 \(\nu\) 表示 G/H 上的测度 \((\kappa\cdot\mu)/\beta\)；由出处同上，它是 G 下拟不变的非零正测度。它的支撑等于 G/H（INT,

VII, 第 10 页, § 1, 第 1 目）。根据 INT, VII, 第 43 页, § 2, 第 2 目, 命题 4，测度 \(\nu\) 是 G/H 上唯一的使得 G 上的测度 \(\nu^{\sharp}\) 等于 \(\kappa \cdot \mu\) 的测度。我们把空间 G/H 赋以测度 \(\nu\)（因而例如我们写 \(\mathcal{L}^{p}(\mathrm{G / H}) = \mathcal{L}^{p}(\mathrm{G / H},\nu)\)）。

 我们称 G 的子集 \(S \subset G\) 是模 H 零测的，如果 \(\varpi(S)\) 是 \(\nu\) -零测的。这一条件不依赖于 G 与 H 上哈尔测度的选取。它蕴含 S 是局部 \(\mu\) -零测的（INT, VII, 第 47 页, § 2, 第 3 目, 命题 6, a)）。我们称 G 的元素所具有的一个性质模 H 几乎处处为真，如果 G 中使该性质不成立的元素所成的集是模 H 零测的。

设 \(\pi\) 是 H 在希尔伯特空间 E 中的一个酉表示。用 \(\mathcal{F}_{\pi}(\mathrm{G})\) 表示 G 上取值于 E 的函数 \(f\) 所成的向量空间，其中 \(f(xh) = \pi(h)f(x)\) 对一切 \((x,h) \in \mathrm{G} \times \mathrm{H}\) 成立。我们称 \(\mathcal{F}_{\pi}(\mathrm{G})\) 的元素为 G 上的 \(\pi\) -等变函数。

G 在 C 上的平凡表示所对应的空间 \(\mathcal{F}_{1}(\mathrm{G})\) 与 G/H 上复值函数的空间 \(\mathcal{F}(\mathrm{G / H})\)，借助映射 \(f\mapsto f\circ \varpi\)（即 \(\mathcal{F}(\mathrm{G / H})\) 到 \(\mathcal{F}_1(\mathrm{G})\) 中的映射）等同。

对任意函数 \(f\)（属于 \(\mathcal{F}_{\pi}(\mathrm{G})\)），函数 \(\|f\|\) 属于 \(\mathcal{F}_{1}(\mathrm{G})\)，因为 \(\pi\) 是酉的。这使我们可以把 \(\|f\|\) 看作 \(\mathrm{G/H}\) 上的一个函数，例如我们将写作 \(\|f(x\mathrm{H})\|\)，以表示该函数在元素 \(x\mathrm{H}\)（\(\mathrm{G/H}\) 的元素）处的值。

设 \(f\) 是 \(\mathcal{F}_{\pi}(\mathrm{G})\) 中的一个函数；若 \(\|f\|\) 在 G/H 的一个紧子集外为零，则称它模 H 在一个紧集外为零。这等价于说，存在 G 的紧子集 K，使 \(f\) 在 \(\mathrm{K} \cdot \mathrm{H}\) 之外为零（TG, III, 第 33 页, 命题 10）。

用 \(\mathcal{K}_{\pi}(\mathrm{G})\) 表示 G 上属于 \(\mathcal{F}_{\pi}(\mathrm{G})\) 且模 H 具有紧支撑的连续函数所成的空间。

与 \(\mathcal{H}_{\pi}(\mathrm{G})\) 类似的一个空间，当 \(\pi\) 是 H 到 \(\mathbf{R}_{+}^{*}\) 中的连续同态时，已在 INT, VII, 第 39 页, § 2, 第 1 目 中定义过。

设 \(p \in [1, +\infty[\)。对每个 \(f \in \mathcal{F}_{\pi}(\mathrm{G})\)，记

\[
\mathrm{N} _ {p} (f) = \left(\int_ {\mathrm{G/H}} ^ {*} \| f \| ^ {p} d \nu\right) ^ {1 / p}.
\]

它是一个实数或 \(+\infty\)。用 \(\mathcal{F}_{\pi}^{p}(\mathrm{G},\nu)\)，或简作 \(\mathcal{F}_{\pi}^{p}(\mathrm{G})\)，表示由函数 \(f\in \mathcal{F}_{\pi}(\mathrm{G})\) 中使 \(\mathrm{N}_p(f)\)

为有限者所成的子空间。空间 \(\mathcal{F}_{\pi}^{p}(G)\) 赋以映射 \(N_{p}\) 后是一个半赋范空间。

空间 \(\mathcal{K}_{\pi}(\mathrm{G})\) 包含于 \(\mathcal{F}_{\pi}^{p}(\mathrm{G})\)。它在 \(\mathcal{F}_{\pi}^{p}(\mathrm{G})\) 中的闭包记作 \(\mathcal{L}_{\pi}^{p}(\mathrm{G},\nu)\)，或简作 \(\mathcal{L}_{\pi}^{p}(\mathrm{G})\)；我们称它为 G 上 \(\pi\) -等变函数中模 H 以指数 \(p^{e}\) 可积的那些函数所成的空间。

下面命题的各断言，当 \(\pi\) 是一维平凡表示时，是 INT, IV, 第 128 页, § 3, 第 3 目, 命题 6 及第 131 页, § 3, 第 4 目, 定理 3 的推论。

命题 3. --- a) 设 \((f_n)_{n \in \mathbf{N}}\) 是 \(\mathcal{F}_\pi^p(\mathrm{G})\) 中的一个序列，使得级数

\[
\sum_ {n = 0} ^ {+ \infty} \mathrm{N} _ {p} (f _ {n})
\]

收敛。于是以 \(f_n(g)\) 为通项的级数对位于一个模 H 零测集 T 之外的每个 \(g\) 绝对收敛。设 \(f\) 是 G 上取值于 E 的函数，它在 G−T 上等于该级数的和，而在 T 上等于 0。这时我们有 \(f \in \mathcal{F}_\pi^p(\mathrm{G})\)，且以 \(f_n\) 为通项的级数收敛于 \(f\)，这里收敛是在空间 \(\mathcal{F}_\pi^p(\mathrm{G})\) 中取的；

\begin{enumerate}
\def\labelenumi{\alph{enumi})}
\setcounter{enumi}{1}
\item
  设 \((f_n)\) 是 \(\mathcal{L}_{\pi}^{p}(\mathrm{G})\) 中收敛于一个函数 \(f\) 的序列。存在 \((f_{n_k})\)——它是 \((f_n)\) 的一个子列——使得 \(f_{n_k}(g)\) 收敛于 \(f(g)\)，对位于一个模 H 零测集之外的每个 \(g\) 成立；
\item
  半赋范空间 \(\mathcal{F}_{\pi}^{p}(\mathrm{G})\) 与 \(\mathcal{L}_{\pi}^{p}(\mathrm{G})\) 是完备的。
\end{enumerate}

在断言 \(a\) ) 中，说以 \(f_n\) 为通项的级数收敛于 \(f\)——这里是指在空间 \(\mathcal{F}_{\pi}^{p}(\mathrm{G})\) 中——意思是部分和的序列 \(f_0 + \cdots + f_n\) 收敛于 \(f\)，收敛是在 \(\mathcal{F}_{\pi}^{p}(\mathrm{G})\) 中取的。这时也称 \(f\) 是该级数的一个和。

我们来证明 \(a\) )。根据 INT, IV, 第 128 页, § 3, 第 3 目, 命题 6，存在一个 \(\nu\) -零测集 S ⊂ G/H，使得以 \(\|f_n(gH)\|\) 为通项的级数对 gH \(\notin\) . S 绝对收敛。此外，对于 gH \(\notin\) . S 等于该级数之和、对于 gH \(\in\) . S 为零的函数 h 满足 N \(_p(h)\) \textless{} + \(\infty\)。

集合 T = \(\overline{\varpi}(S)\) 是模 H 零测的。对每个 \(g \notin T\)，以 \(f_{n}(g)\) 为通项的级数在 E 中绝对收敛。令 \(f(g) = \sum f_{n}(g)\)（对 \(g \notin T\)），否则令 \(f(g) = 0\)。我们有 \(f \in \mathcal{F}_{\pi}(G)\)。注意 \(\|f(gH)\| \leqslant h(gH)\) 对所有 \(g \in G\) 成立，由此 \(N_{p}(f) \leqslant\)

\(N_{p}(h) < +\infty\)。于是我们有 \(f \in \mathcal{F}_{\pi}^{p}(G)\)。类似地，由此得出

\[
\mathrm{N} _ {p} \Big (f - \sum_ {n = 0} ^ {k} f _ {n} \Big) \leqslant \sum_ {n = k + 1} ^ {+ \infty} \mathrm{N} _ {p} (f _ {n})
\]

对所有 \(k \in \mathbf{N}\) 成立，因此级数 \(\sum f_n\) 收敛于 \(f\)，收敛是在 \(\mathcal{F}_{\pi}^{p}(\mathrm{G})\) 中取的。断言 \(a)\) 得证。

我们来证明 \(b\) )。序列 \((\|f_n - f\|)_{n \in \mathbf{N}}\) 在空间 \(\mathcal{L}^p(\mathrm{G}/\mathrm{H})\) 中收敛于 0。根据 INT, IV, 第 131 页, § 4, 第 3 目, 定理 3，存在一个子序列 \((\|f_{n_k} - f\|)_{k \in \mathbf{N}}\)，它 \(\nu\) -几乎处处收敛于 0。于是序列 \((f_{n_k}(g))_{k \in \mathbf{N}}\) 收敛于 \(f(g)\)，对位于一个模 H 零测集之外的每个 \(g\) 成立。

最后证明 \(c\) )。设 \((f_n)_{n \in \mathbf{N}}\) 是 \(\mathcal{F}_\pi^p(\mathrm{G})\) 中的一个柯西序列。存在严格递增的整数序列 \((n_k)_{k \in \mathbf{N}}\)，使得 \(\mathrm{N}_p(f_{n_{k+1}} - f_{n_k}) \leqslant 2^{-k}\) 对所有 \(k \in \mathbf{N}\) 成立。对每个 \(k \in \mathbf{N}\)，令 \(h_k = f_{n_{k+1}} - f_{n_k} \in \mathcal{F}_\pi^p(\mathrm{G})\)。由 \(a\) )，以 \(h_k\) 为通项的级数在 \(\mathcal{F}_\pi^p(\mathrm{G})\) 中收敛；记其和为 \(h\)。对每个 \(\ell \in \mathbf{N}\)，由此得出

\[
f _ {n _ {0}} + \sum_ {k = 0} ^ {\ell} h _ {k} = f _ {n _ {\ell + 1}},
\]

于是 \(f_{n_{0}} + h\) 是序列 \((f_{n})\) 的一个聚点。从而该序列收敛（TG, II, 第 14 页, 命题 5 的推论 2）。这样，空间 \(\mathcal{F}_{\pi}^{p}(\mathrm{G})\) 是完备的；又因空间 \(\mathcal{L}_{\pi}^{p}(\mathrm{G})\) 在 \(\mathcal{F}_{\pi}^{p}(\mathrm{G})\) 中是闭的，它也是完备的。

推论. --- 设 \((f_{n})_{n\in\mathbf{N}}\) 是 \(\mathcal{L}_{\pi}^{p}(\mathrm{G})\) 中的一个柯西序列，并设 \(f\in\mathcal{F}_{\pi}(\mathrm{G})\) 使得 \(f_{n}(g)\) 模 H 几乎处处收敛于 \(f(g)\)。则 \(f\in\mathcal{L}_{\pi}^{p}(\mathrm{G})\)，且 \((f_{n})_{n\in\mathbf{N}}\) 在 \(\mathcal{L}_{\pi}^{p}(\mathrm{G})\) 中收敛于 f。

事实上，函数 \(f\) 与序列 \((f_n)\) 在 \(\mathcal{L}_{\pi}^{p}(\mathrm{G})\) 中的极限模 H 几乎处处重合。

用 \(\mathrm{L}_{\pi}^{p}(\mathrm{G})\) 表示与半赋范空间 \(\mathcal{L}_{\pi}^{p}(\mathrm{G})\) 相伴随的分离赋范拓扑向量空间；它是一个巴拿赫空间。

设 \(f\) 是 \(\mathrm{G}\) 上取值于 \(\mathrm{E}\) 的一个函数。用 \(\mathrm{S}_{f}\) 表示由 \(x\in \mathrm{G}\) 中使函数 \(h\mapsto f(xh)\)（定义在 \(\mathrm{H}\) 上）不属于 \(\mathcal{L}_{\mathrm{E}}^{1}(\mathrm{H})\) 者所成的集。我们有 \(\mathrm{S}_f\cdot h = \mathrm{S}_f\) 对所有 \(h\in \mathrm{H}\) 成立。

设 \(x \in \mathrm{G} - \mathrm{S}_f\)。由于 \(\pi\) 是酉表示，函数 \(h \mapsto \pi(h)^* f(xh)\) 可测（它是 H 上的映射 \(h \mapsto (h, f(xh))\)——后者到 \(\mathrm{H} \times \mathrm{E}\) 中可测——与 H × E 到 E 中的连续映射 \((g,x)\mapsto\pi(g)^{*}x\) 的复合，参见 INT, IV, 第 174 页, § 5, 第 3 目, 定理 1），并且它在 H 上可积。

我们通过下式在 G 上定义一个函数 \(f^{\pi}\)：

\[
f ^ {\pi} (x) = \int_ {\mathrm{H}} \pi (h) ^ {*} f (x h) d \beta (h)\tag{3}
\]

其中 \(x \in \mathrm{G} - \mathrm{S}_f\)；而 \(f^\pi(x) = 0\)（当 \(x \in \mathrm{S}_f\) 时）。

命题 4. --- 设 \(f \in \mathcal{L}_{\mathrm{E}}^{1}(\mathrm{G})\)。

\begin{enumerate}
\def\labelenumi{\alph{enumi})}
\item
  集合 \(\mathrm{S}_f\) 模 H 零测，且 \(f^\pi \in \mathcal{F}_\pi(\mathrm{G})\)；
\item
  设 C 是 G 的一个紧子集。若 f 连续且支撑含于 C，则 \(S_{f}\) 为空，函数 \(f^{\pi}\) 属于 \(\mathcal{H}_{\pi}(\mathrm{G})\) 且其支撑含于 \(C \cdot H\)。
\end{enumerate}

\(a\) ) 的第一部分由 INT, VII, 第 57 页, § 2, 第 5 目, c) 得出。设 \(w \in H\)。若 \(x \in G - S_f\)，则 \(xw \in G - S_f\)，且

\[
\begin{array}{c} f ^ {\pi} (x w) = \int_ {\mathrm{H}} \pi (h) ^ {*} f (x w h) d \beta (h) \\ = \int_ {\mathrm{H}} \pi (w ^ {- 1} y) ^ {*} f (x y) d \beta (y) = \pi (w) f ^ {\pi} (x), \end{array}
\]

而若 \(x \in S_{f}\)，则 \(f^{\pi}(xw) = 0 = \pi(w)^{*}f^{\pi}(x)\)。因此函数 \(f^{\pi}\) 属于 \(\mathcal{F}_{\pi}(\mathrm{G})\)。

现在设 \(f\) 连续且支撑含于 C。于是对每个 \(x \in G\)，映射 \(h \mapsto f(xh)\) 属于 \(\mathcal{K}(H)\)，从而可积，这就证明了 \(S_f = \varnothing\)。

我们来证明 \(f^{\pi}\) 连续。设 \(x \in G\)。设 U 是 \(x\) 在 G 中一个相对紧的开邻域。

对每个 \(y \in \mathrm{U}\)，我们有

\[
\begin{array}{l} \| f ^ {\pi} (y) - f ^ {\pi} (x) \| \leqslant \int_ {\mathrm{H}} \| f (y h) - f (x h) \|   d \beta (h) \\ \qquad = \int_ {\mathrm{H} \cap (y ^ {- 1} \mathrm{C} \cup x ^ {- 1} \mathrm{C})} \| f (y h) - f (x h) \|   d \beta (h) \\ \qquad \leqslant \beta (\mathrm{H} \cap \mathrm{U} ^ {- 1} \mathrm{C}) \sup _ {h \in \mathrm{U} ^ {- 1} \mathrm{C}} \| f (y h) - f (x h) \|, \end{array}
\]

而 \(f^{\pi}\) 的连续性则由 f 在 G 上的一致连续性得出。

最后，若 \(x \in G\) 满足 \(f^{\pi}(x) \neq 0\)，则存在 \(h \in H\) 使 \(f(xh) \neq 0\)，于是 xh 属于 C，而 x 属于 C·H。由于 C·H 在 G 中闭，我们断定 \(f^{\pi}\) 的支撑含于 C·H。

 设 C 是 G 的一个紧子集。设 \(u \in \mathcal{K}_{+}(G)\) 是满足 \(u(x) > 0\)（对一切 \(x \in C\)）的函数。设 v 是 G 上由下式定义的函数

\[
v (x) = \int_ {\mathrm{H}} u (x h) d \beta (h)\tag{4}
\]

对所有 \(x \in G\) 成立；沿用前面的记号，我们有 \(v = u^{1}\)，它对应于 H 的一维平凡表示。函数 v 连续且为正；它属于 \(\mathcal{F}_{1}(G)\)，其支撑含于 \(\operatorname{Supp}(u) \cdot H\)，并且我们有

\[
\inf _ {x \in \mathrm{C} \cdot \mathrm{H}} v (x) > 0\tag{5}
\]

（INT, VII, 第 39--40 页, § 2, 第 1 目, 命题 1 及引理 1, a)）。

引理 5. --- 设 \(f \in \mathcal{F}_{\pi}(\mathrm{G})\) 是一个 \(\mu\) -可测且在 \(\mathrm{C} \cdot \mathrm{H}\) 之外为零的函数，使得函数 \(\|f\|\) \(\nu\) -可积（在 \(\mathrm{G}/\mathrm{H}\) 上）。设 s 是 \(\mathrm{G}\) 上取值于 E 的函数，满足

\[
 s (x) = \left\{ \begin{array}{l l} v (x) ^ {- 1} f (x) & \text{若 } x \in \mathrm{C} \cdot \mathrm{H} \\ 0 & \text{若 } x \in \mathrm{G} - \mathrm{C} \cdot \mathrm{H}. \end{array} \right.
\]

\begin{enumerate}
\def\labelenumi{\alph{enumi})}
\item
  函数 s 是 \(\mu\) -可测的；它属于 \(\mathcal{F}_{\pi}(\mathrm{G})\) 且在 C·H 之外为零；
\item
  函数 us 属于 \(\mathcal{L}^1 (\mathrm{G})\)，且模 H 几乎处处有 \((us)^{\pi} = f\)。
\end{enumerate}

函数 \(s\) 按定义在 C \(\cdot\) . H 之外为零；它是 \(\mu\) -可测的（INT, IV, 第 193 页, § 5, 第 10 目, 命题 16），因为函数 \(f\) \(\mu\) -可测且 \(v(x) > 0\) 对一切 \(x \in \mathrm{C} \cdot \mathrm{H}\) 成立。我们有 \(s \in \mathcal{F}_{\pi}(\mathrm{G})\)，因为 \(v \in \mathcal{F}_1(\mathrm{G})\)。

函数 \(f\) 是局部 \(\mu\) -可积的，因为 \(\|f\|\) 是 \(\nu\) -可积的函数（INT, VII, 第 47 页, § 2, 第 3 目, 命题 6, c)，注意测度 \(\nu^{\sharp} = \kappa \cdot \mu\) 等价于 \(\mu\)），从而由公式 (5)，函数 \(s\) 亦然。函数 \(us\) 可测且在 G 的一个紧子集之外为零；由于 \(u\) 有界，由此得出 \(us\) 是 \(\mu\) -可积的。

对每个 \(x\in \mathrm{G - S}_{us}\)，我们有

\[
\begin{array}{c} (u s) ^ {\pi} (x) = \int_ {\mathrm{H}} \pi (h) ^ {*} (u (x h) s (x h)) d \beta (h) \\ = \int_ {\mathrm{H}} u (x h) s (x) d \beta (h) = v (x) s (x), \end{array}
\]

因为 \(s(xh) = \pi(h)s(x)\)；最后一个断言由命题 4, a) 得出。

当 \(H = \{e\}\) 且 \(\pi\) 是一维时，下面的命题就是 INT, IV, 第 184 页, § 5, 第 6 目, 定理 5。

命题 5. --- 设 \(p \in [1, +\infty[\)。空间 \(\mathcal{L}_{\pi}^{p}(\mathrm{G})\) 是由满足如下条件的函数 \(f \in \mathcal{F}_{\pi}(\mathrm{G})\) 所成的空间：\(f\) \(\mu\) -可测，且函数 \(\|f\|\) 属于 \(\mathcal{L}^{p}(\mathrm{G}/\mathrm{H})\)。

设 \(f \in \mathcal{L}_{\pi}^{p}(G)\)。它是 \(\mathcal{L}_{\pi}^{p}(G)\) 中由 \(\mathcal{K}_{\pi}(G)\) 的元素组成的序列的极限。由命题 3, b) 及叶戈罗夫定理（INT, IV, 第 175 页, § 5, 第 4 目, 定理 2），函数 \(f\) 因此是 \(\mu\) -可测的；从而 G/H 上的函数 \(\|f\|\) 是 \(\nu\) -可测的（INT, VII, 第 47 页, § 2, 第 3 目, 命题 6, b)）。由于 \(N_{p}(f)\) 有限，函数 \(\|f\|\) 属于 \(\mathcal{L}^{p}(G/H)\)（INT, IV, 第 184 页, § 5, 第 6 目, 定理 5）。

 我们来证明逆命题。设 \(f\) 是一个 \(\mu\) -可测、属于 \(\mathcal{F}_{\pi}(\mathrm{G})\) 的函数，且满足 \(\|f\|\in\mathcal{L}^{p}(\mathrm{G}/\mathrm{H})\)。设 \(\varepsilon>0\)。我们来证明存在 \(\widetilde{f}\in\mathcal{K}_{\pi}(\mathrm{G})\)，使得

\[
\mathrm{N} _ {p} (f - \widetilde {f}) ^ {p} = \int_ {\mathrm{G/H}} ^ {*} \| f - \widetilde {f} \| ^ {p} d \nu \leqslant \varepsilon ,
\]

如此即完成证明。

 先设存在 G 的紧子集 C，使 \(f\) 在 \(\mathrm{C} \cdot \mathrm{H}\) 之外为零。设 \(u \in \mathcal{K}_{+}(\mathrm{G})\) 是满足 \(u(x) > 0\)（对一切 \(x \in \mathrm{C}\)）的函数，并按上式定义 \(v = u^{1}\)。用 \(\varphi\) 表示 \(u\) 的支撑的特征函数。

设 q 是 p 的共轭指数。若 p = 1，设 w 是 G 上等于 \(\sup_{x \in G} u(x)\) 的常值函数。若 p \textgreater{} 1，我们定义

\[
w (x) = \left(\int_ {\mathrm{H}} u (x h) ^ {q} d \beta (h)\right) ^ {1 / q}
\]

 对所有 \(x \in G\) 成立。在所有情形下，函数 w 都连续且为正；它属于 \(\mathcal{K}_{1}(G)\)（命题 4 应用于一维平凡表示及函数 \(u^{q}\)），从而在 G 上有界。我们令 \(W = \sup_{x \in G} w(x)\)。

 设 \(s\) 是 \(\mathbf{G}\) 上取值于 \(\mathbf{E}\) 的、把引理 5 应用于 \(f\) 时所得到的函数。令 \(g = \varphi s\)。函数 \(g\) 是 \(\mu\) -可测的，且满足 \(\| g \| \leqslant \| f \| / \inf_{x \in \mathrm{C} \cdot \mathrm{H}} v(x)\)。由于 \(g\) 在 \(u\) 的支撑之外为零且 \(\kappa\) 连续，我们有 \(g \in \mathcal{L}_{\mathrm{E}}^{p}(\mathrm{G}, \kappa \cdot \mu)\)。设 \(\widetilde{g}\) 是

\(\mathcal{K}_{\mathrm{E}}(\mathrm{G})\) 中的一个函数，使得

\[
\int_ {\mathrm{G}} ^ {*} \| g - \widetilde {g} \| ^ {p} \kappa d \mu \leqslant \frac {\varepsilon}{\mathrm{W} ^ {p}}.
\]

我们有 us = ug，从而模 H 几乎处处有 \(f = (us)^{\pi} = (ug)^{\pi}\)（引理 5, b)）。令 \(\widetilde{f} = (u\widetilde{g})^{\pi}\)；我们有 \(\widetilde{f} \in \mathcal{K}_{\pi}(\mathrm{G})\)（命题 4, b)）。对每个 \(x \in G - S_{ug}\)，我们得到

\[
\| (u g) ^ {\pi} (x) - (u \widetilde {g}) ^ {\pi} (x) \| \leqslant \int_ {\mathrm{H}} ^ {*} u (x h) \| g (x h) - \widetilde {g} (x h) \| d \beta (h)
\]

由此得

\[
\| (u g) ^ {\pi} (x) - (u \widetilde {g}) ^ {\pi} (x) \| ^ {p} \leqslant w (x) ^ {p} \int_ {\mathrm{H}} ^ {*} \| g (x h) - \widetilde {g} (x h) \| ^ {p} d \beta (h),
\]

这里在 \( p \) \textgreater{} 1 的情形用了赫尔德不等式。由于 \(S_{ug}\) 模 H 零测，由此得出

\[
\begin{array}{c} \int_ {\mathrm{G/H}} ^ {*} \| f - \widetilde {f} \| ^ {p} d \nu \leqslant \mathrm{W} ^ {p} \int_ {\mathrm{G/H}} ^ {*} \Big (\int_ {\mathrm{H}} ^ {*} \| g (x h) - \widetilde {g} (x h) \| d \beta (h) \Big) ^ {p} d \nu (x \mathrm{H}) \\ = \mathrm{W} ^ {p} \int_ {\mathrm{G}} ^ {*} \| g - \widetilde {g} \| ^ {p} \kappa d \mu \leqslant \varepsilon \end{array}
\]

这是根据 INT, VII, 第 46 页, § 2, 第 3 目, 命题 5, b)，它适用于 \(g - \widetilde{g}\) 在 G 的一个紧子集之外为零的情形。这就证明了当 \(f\) 模 H 在一个紧子集之外为零时所要求的性质。

现在考虑一般情形。由假设 \(\|f\|\in\mathcal{L}^{p}(G/H)\)，存在 G/H 的紧子集 L，使得

\[
\int_ {(\mathrm{G} / \mathrm{H}) - \mathrm{L}} ^ {*} \| f \| ^ {p} d \nu \leqslant \frac {\varepsilon}{2}
\]

（参见 INT, IV, 第 152 页, § 4, 第 6 目, 定理 4）。设 \(\varphi_{\mathrm{L}}\) 是 L 的特征函数，并令 \(f_{\mathrm{L}} = (\varphi_{\mathrm{L}} \circ \varpi)f\)。我们有

\[
\int_ {\mathrm{G} / \mathrm{H}} ^ {*} \| f - f _ {\mathrm{L}} \| ^ {p} d \nu = \int_ {(\mathrm{G} / \mathrm{H}) - \mathrm{L}} ^ {*} \| f \| ^ {p} d \nu \leqslant \frac {\varepsilon}{2}.
\]

函数 \(f_{\mathrm{L}}\) 是 \(\mu\) -可测的且模 H 在一个紧集之外为零。它属于 \(\mathcal{L}_{\pi}^{p}(\mathrm{G})\)，于是由前一种情形，存在 \(\widetilde{f} \in \mathcal{K}_{\pi}(\mathrm{G})\) 使得 \(\mathrm{N}_p(f_{\mathrm{L}} - \widetilde{f}) \leqslant (\frac{\varepsilon}{2})^{1/p}\)，从而

\[
\int_ {\mathrm{G/H}} ^ {*} \| f - \widetilde {f} \| ^ {p} d \nu \leqslant \varepsilon ,
\]

这正是所要证的。

 我们来考虑 \(p = 2\) 的情形。对 \(f_1\) 与 \(f_2\)（均属于 \(\mathcal{F}_{\pi}(\mathrm{G})\)），函数 \(x \mapsto \langle f_1(x) | f_2(x) \rangle\) 属于 \(\mathcal{F}_1(\mathrm{G})\)，因为表示 \(\pi\) 是酉的，从而通过取商定义了 \(\mathrm{G / H}\) 上的一个函数，我们和先前一样把它等同于 \(\langle f_1 | f_2 \rangle\)。我们有不等式 \(|\langle f_1 | f_2 \rangle| \leqslant \| f_1 \| \| f_2 \|\)（在 \(\mathcal{F}(\mathrm{G / H})\) 中）。

若 \(f_1\) 与 \(f_2\) 属于 \(\mathcal{L}_\pi^2(\mathrm{G})\)，则函数 \(\langle f_1 | f_2 \rangle\) 属于 \(\mathcal{L}_1^1(\mathrm{G})\)。特别地，它在 \(\mathrm{G/H}\) 上可积。映射

\[
(f _ {1}, f _ {2}) \mapsto \int_ {\mathrm{G} / \mathrm{H}} \langle f _ {1} | f _ {2} \rangle d \nu
\]

是 \(\mathcal{L}_{\pi}^{2}(\mathrm{G})\) 上的一个正埃尔米特形式；其相应的半范数是半范数 \(N_{2}\)。特别地，空间 \(\mathcal{L}_{\pi}^{2}(\mathrm{G})\) 是一个预希尔伯特空间，而 \(\mathrm{L}_{\pi}^{2}(\mathrm{G})\) 是与 \(\mathcal{L}_{\pi}^{2}(\mathrm{G})\) 相伴随的希尔伯特空间。

